%% 联合培养企业评价意见正文
%% 适用学位：专项专硕（SpecialEnMaster）、实践总结报告（PracticeReport）
[请重点对以下两方面进行评价：

一是学位论文的实用性、先进性、新颖性，包括结合实践进行的理论创新、工业指标上的突破、新方案的提出、新装备的生产、新装置的设计等，重点关注学位论文相关成果在实践应用中的落地情况与贡献程度；

二是作者工程能力，可从技术能力、问题解决能力、项目管理能力和团队合作能力等方面进行综合评价。

以下为示例。]

XX同学的学位论文工作，依托企业/校企联合的XX项目/XX关键任务开展，经过深入研究XX相关理论，结合实际情况提出了具有建设性的见解和解决方案，取得了XX成果，为企业突破XX工程瓶颈问题/技术创新/产业发展提供了有力支撑。

XX同学在企业实践期间，充分展现了其坚实的理论基础和系统的专业知识，以及理论与实践相结合解决实际问题的能力，同时表现出良好的职业道德、职业素养和出色的团队合作精神。

他/她对项目的研究过程和成果进行了详尽、系统的总结和提炼，依托企业实践期间取得的成果，完成了本学位论文，达到了申请硕士学位的标准。同意XX同学以本学位论文申请答辩。

\vspace{6ex}

\hfill\begin{minipage}[t]{0.5\textwidth}
联合培养企业：

盖章日期：\hspace{4\ccwd}年\hspace{2\ccwd}月\hspace{2\ccwd}日
\end{minipage}
